% !TeX root = ../main.tex

\section{Resultados e análise}
\label{sec:resultados}

\subsection{Protocolo experimental}
\preencher{Caracterizar as instâncias, a origem e a preparação dos dados,
o ambiente computacional e os cenários comparados. Definir as métricas e
o cronograma de referência, garantindo condições comparáveis entre os
experimentos.}

\subsection{Resultados computacionais}
\preencher{Apresentar somente resultados obtidos e verificados. Informar,
por instância e cenário, prazo, tempo de execução, estado da solução e gap
quando disponível. Numerar tabelas e figuras e mencioná-las no texto.
Distinguir solução viável de ótimo comprovado.}

\subsection{Discussão e limitações}
\preencher{Interpretar os ganhos e gargalos observados e relacioná-los às
hipóteses e aos trabalhos discutidos. Explicar as limitações dos dados,
da representação dos recursos e do método, delimitando o alcance das
conclusões.}
